% Version fixed135: 9- and 27-sheet odd-index tower, actual Sylow-2
% preimage for Q2(2^(1/9),zeta_9), exact marked Fox matrices,
% integral pro-2 primitive Schreier/Tietze criterion,
% and nested chain-level Hecke idempotents of ranks 1,2,6,18.
% Integrated into the fixed135 complete mother TeX.

\subsection{An odd-index tame tower beyond $S_3$:
nine sheets, twenty-seven Sylow sheets and coherent Fox--Hecke descent}
\label{subsec:dyadic-f135-27sheet-sylow-tower}

The direct three-sheet calculation of fixed133--134
has a genuine odd-index extension.
The degree-nine Eisenstein field is not itself the
Sylow-$2$ preimage of its Galois closure:
the residue stabilizer contains a nontrivial
odd tame factor.  The correct Sylow field has
degree twenty-seven.
We construct \emph{both} covers from the actual
marked full-group relations, give their complete
integral coinduced Fox differentials and
elementary divisors, reduce the marked
Schreier group relators to actual one-relator
pro-$2$ presentations, and construct a
strictly nested four-stage Hecke projector.
This is an explicit enlargement from the
single $S_3$ sector to a different tame
residual quotient of order fifty-four.

Choose $\alpha^9=2$ and put
\begin{equation}\label{eq:dyadic-f135-fields}
\boxed{
\begin{gathered}
F_3=\mathbf Q_2(\alpha^3),\qquad
F_9=\mathbf Q_2(\alpha),\\
L_9=F_9(\zeta_9),\qquad
K_{27}=L_9^{\langle\sigma^3\rangle}.
\end{gathered}}
\end{equation}
Choose the tame marked generators
$\sigma,\tau$ so that the finite quotient
on $L_9$ satisfies
\begin{equation}\label{eq:dyadic-f135-tame-54-quotient}
\boxed{
Q_9=\operatorname{Gal}(L_9/\mathbf Q_2)
\simeq C_9\rtimes C_6,\qquad
\tau^9=\sigma^6=1,\quad
\sigma^{-1}\tau\sigma=\tau^2.
}
\end{equation}
In this quotient the marked wild
generators $x_0,x_1$ have trivial image.
Let
\[
\Gamma=G_{\mathbf Q_2},\quad
H_3=G_{F_3},\quad
H_9=G_{F_9},\quad
H_{27}=G_{K_{27}}.
\]

\begin{theorem}[The actual Sylow-$2$ field and three index-three steps]
\label{thm:dyadic-f135-actual-27-sylow-field}
The finite extensions in
\eqref{eq:dyadic-f135-fields}
have
\begin{equation}\label{eq:dyadic-f135-odd-field-degrees}
[\Gamma:H_3]=3,\qquad
[H_3:H_9]=3,\qquad
[H_9:H_{27}]=3.
\end{equation}
The full tame quotient $Q_9$
has order $54$, and its
Sylow $2$-subgroup is
$\langle\sigma^3\rangle\simeq C_2$.
Its inverse image in $\Gamma$
is precisely $H_{27}$:
\begin{equation}\label{eq:dyadic-f135-sylow-preimage-27}
\boxed{\quad
[\,\Gamma:H_{27}\,]=27,\qquad
\pi(H_{27})=\langle\sigma^3\rangle.
\quad}
\end{equation}
More concretely, if $U_3$
denotes the unique unramified
cubic extension of $\mathbf Q_2$
inside $\mathbf Q_2(\zeta_9)$,
then
\[
K_{27}=F_9U_3,
\qquad [K_{27}:\mathbf Q_2]=27.
\]
Each of $F_9$ and $K_{27}$
is a finite dyadic field of
odd degree and contains
exactly $\mu_2$ among
its $2$-power roots of unity.
Its maximal pro-$2$
Galois quotient is a
Demu\v{s}kin group, with
\begin{equation}\label{eq:dyadic-f135-demushkin-odd-ranks}
\boxed{
d(G_{F_9}(2))=11,\qquad
d(G_{K_{27}}(2))=29,\qquad
q=2
}
\end{equation}
and full cyclotomic orientation
image $\mathbf Z_2^\times$.
The degree-nine field
$F_9$ is \emph{not}
itself a Sylow-$2$
preimage for $Q_9$:
its image in $Q_9$
is $C_6$.
\end{theorem}

\begin{proof}
The polynomial $X^9-2$
is Eisenstein at $2$,
so $F_9/\mathbf Q_2$
is totally and tamely
ramified of degree nine;
its cubic subfield is
$F_3$.
Since $9$ is odd,
$\mathbf Q_2(\zeta_9)$
is unramified of degree
$\operatorname{ord}_9(2)=6$.
The two fields are linearly
disjoint.
The tame inertia action on
$\alpha$ and the Frobenius
action on $\zeta_9$ give
the displayed semidirect
quotient, in the chosen
inverse-Frobenius convention
for $\sigma$.
Its $2$-Sylow has order
two and is generated by
$\sigma^3$.
The stabilizer of the
chosen ninth root
$\alpha$ is
$\langle\sigma\rangle
\simeq C_6$;
the stabilizer of
$\alpha^3$ is
$\langle\tau^3,\sigma\rangle$
of order eighteen.
The intermediate
subgroup inclusions
therefore have index
three, as claimed.

The subgroup
$\langle\sigma^3\rangle$
fixes both $\alpha$
and the unramified
cubic subfield of
$\mathbf Q_2(\zeta_9)$.
The compositum of
those fixed fields
has degree $9\cdot3=27$,
hence is exactly the
fixed field of this
order-two subgroup.
All field degrees
are odd, so no field
in the tower contains
$\mu_4$.
The cyclotomic pro-$2$
extension has trivial
intersection with an
odd-degree field;
the cyclotomic
orientation remains
surjective onto
$\mathbf Z_2^\times$.
The local Demu\v{s}kin
rank formula
$d=[K:\mathbf Q_2]+2$
then gives $11,29$
and the exact $q=2$
power invariant.
\end{proof}

\begin{theorem}[Explicit nine- and twenty-seven-coset tame permutations]
\label{thm:dyadic-f135-permutation-9-27}
Let $V_n=\mathbf Z_2^n$
be the permutation
module on the chosen
left cosets
$H_n\backslash\Gamma$,
where $n=9$ or $27$.
For $n=9$ use the
representatives
$t_i=\tau^i$,
$i\in\mathbf Z/9$,
and denote their basis
vectors by $e_i$.
The coefficient
actions are
\begin{equation}\label{eq:dyadic-f135-nine-permutations}
\boxed{
\Sigma_9 e_i=e_{5i},
\qquad
\Upsilon_9e_i=e_{i-1}
\quad(i\in\mathbf Z/9).
}
\end{equation}

For $n=27$ use
the representatives
\begin{equation}\label{eq:dyadic-f135-27-coset-representatives}
\boxed{\qquad
t_{i,j}=\tau^i\sigma^j,\qquad
i\in\mathbf Z/9,\quad j=0,1,2.
\qquad}
\end{equation}
Their basis vectors
$e_{i,j}$ have
the explicit
coefficient actions
\begin{equation}\label{eq:dyadic-f135-27-permutations}
\boxed{
\begin{aligned}
\Upsilon_{27}e_{i,j}
 &=e_{i-5^j,j},\\
\Sigma_{27}e_{i,j}
 &=
 \begin{cases}
 e_{i,j-1},&j=1,2,\\
 e_{-i,2},&j=0.
 \end{cases}
\end{aligned}}
\end{equation}
All indices $i$ are
taken modulo nine.
Both pairs satisfy
\begin{equation}\label{eq:dyadic-f135-permutation-tame-relation}
\boxed{
\Sigma_n^{-1}\Upsilon_n\Sigma_n
=\Upsilon_n^2,\qquad
\Sigma_n^6=I,\qquad
\Upsilon_n^9=I.
}
\end{equation}
The $\Upsilon_9$ action
has one orbit of length nine,
and $\Upsilon_{27}$
has three orbits of
length nine, cyclically
permuted by $\Sigma_{27}$.
Thus the integral
tame-inertia projector
\begin{equation}\label{eq:dyadic-f135-odd-inertia-idempotent}
\boxed{\qquad
P_n=\frac19\sum_{r=0}^8\Upsilon_n^r,
\qquad P_n^2=P_n
\qquad}
\end{equation}
has respective ranks
\[
\operatorname{rank}P_9=1,\qquad
\operatorname{rank}P_{27}=3.
\]
\end{theorem}

\begin{proof}
For $n=9$ the chosen
right-coset multiplication
is
\[
\tau^i\tau=\tau^{i+1},
\qquad
\tau^i\sigma=\sigma\tau^{2i}.
\]
Coefficient coinduction
uses the inverse
right-coset permutation,
which gives
\eqref{eq:dyadic-f135-nine-permutations}.
For $n=27$, write
$S=\langle\sigma^3\rangle$
and choose the representatives
$\tau^i\sigma^j$ with
$j=0,1,2$.
In $Q_9$ the relation
$\sigma\tau\sigma^{-1}
=\tau^5$ gives
\[
(\tau^i\sigma^j)\tau
 =\tau^{i+5^j}\sigma^j.
\]
Right multiplication
by $\sigma$ increments
$j$ until $j=2$;
the relation
$\sigma^3\tau^{-i}
=\tau^i\sigma^3$
then sends
$(i,2)$ to
$(-i,0)$.
Taking the inverse
permutation proves
\eqref{eq:dyadic-f135-27-permutations}.
The displayed
tame relation follows
by substitution.
Each $\tau$ orbit
has length nine.
For $n=27$
the three orbits
are indexed by $j$,
and the $\sigma$
action cycles them.
Since each orbit
has odd length,
the normalized norm
$P_n$ is an
integral idempotent
whose image consists
of vectors constant
on each orbit.
\end{proof}

\begin{theorem}[Uniform nine-/twenty-seven-sheet marked Fox differential]
\label{thm:dyadic-f135-27sheet-fox-block-formula}
Let
$\mathcal J_n^\bullet$
be the literal
two-marked-relator
cochain complex of
$G_{\mathbf Q_2}$
with coefficient
module $V_n$,
$n\in\{9,27\}$:
\[
\mathcal J_n^\bullet:
\quad
V_n\xrightarrow{d^0_n}
V_n^{\oplus4}
\xrightarrow{d^1_n}
V_n^{\oplus2}.
\]
The generator order is
$(\sigma,\tau,x_0,x_1)$;
the relation order is
(tame,wild).
Then the \emph{complete}
integral differentials
are
\begin{equation}\label{eq:dyadic-f135-27sheet-fox-matrix}
\boxed{
\begin{aligned}
d_n^0(v)&=
\bigl((\Sigma_n-I)v,\,
      (\Upsilon_n-I)v,\,
      0,0\bigr),\\
d_n^1(a,b,c,d)
&=\left(
 \Sigma_n^{-1}(\Upsilon_n-I)a
 +(\Sigma_n^{-1}-I-\Upsilon_n)b,\right.\\
&\hspace{3.7em}\left.
 3P_nb+(4P_n-2I)c
 +(\Sigma_n^{-1}-P_n)d
\right).
\end{aligned}}
\end{equation}
All entries belong to
$\mathbf Z_2$:
the \emph{only}
noninteger coefficient
is division by
the odd unit nine.
The displayed matrices
satisfy the literal
chain identity
\begin{equation}\label{eq:dyadic-f135-27sheet-dsq}
\boxed{\qquad
d_n^1d_n^0=0.
\qquad}
\end{equation}
They have degrees
\begin{equation}\label{eq:dyadic-f135-n-sheet-complex-degrees}
\boxed{
(n,4n,2n)=
\begin{cases}
(9,36,18),\\
(27,108,54).
\end{cases}}
\end{equation}

The same formula
applies to
\emph{any}
odd-index transitive
permutation quotient
of the chosen
marked $G_{\mathbf Q_2}$
source for which
$\Upsilon$ has odd
order and the
pro-$2$ component
of $\Sigma$ has
linear square one;
the $P$ in that
formula is always
the normalized
odd-inertia average.
The present
two concrete
permutation modules
satisfy these
hypotheses.
\end{theorem}

\begin{proof}
Use crossed-derivation
coordinates in the
affine semidirect
extension
$V_n\rtimes
\langle\Sigma_n,\Upsilon_n\rangle$.
The wild generators
have trivial
linear action,
so $x_0,x_1$
are pure translations.
The residual
$\tau$ permutation
has order nine;
therefore the
pro-$2$ projection
of $(x_j\tau)$
is the pure
translation
\[
u_j=P_n(b+c_j),
\]
where $b$ is the
translation at
$\tau$ and $c_j$
at $x_j$.
Indeed its ninth
power has
translation
$(I+\Upsilon+\cdots+
\Upsilon^8)(b+c_j)$,
and the unique
pro-$2$ ninth root
is exactly $P_n(b+c_j)$.

The residual
$\sigma$ permutation
has order six.
Its $2$-primary
component $\sigma_2$
has residual
linear part
$\Sigma_n^3$
of order two.
Hence
$\sigma_2^2$
acts by a pure
translation and
centralizes all
the wild
translations.
As in the
fixed128--129
affine word
calculation, the
literal Roe--Turturean
wild auxiliary
words simplify to
\[
d_0=P_n(b+c)-c,
\qquad
h_0=2c+4d_0,
\qquad
u_1=P_n(b+d),
\]
while the
remaining commutators
of pure translations
vanish.
The conjugated
$x_1^\sigma$
contributes
$\Sigma_n^{-1}d$.
Thus the
wild relation has
translation
\[
h_0-u_1+\Sigma_n^{-1}d
=3P_nb+(4P_n-2I)c+
(\Sigma_n^{-1}-P_n)d.
\]
The tame relation
$\sigma^{-1}\tau\sigma
\tau^{-2}=1$
has its standard
Fox differential
\[
\Sigma_n^{-1}(\Upsilon_n-I)a+
(\Sigma_n^{-1}-I-\Upsilon_n)b.
\]
The principal
crossed derivation
$d^0$ is
$(\rho(g)-I)v$
on each marked
generator, giving
the first line.
The two relation
matrices evaluate
to identity at
the source
permutation
representation,
so Fox's fundamental
formula gives
$d^1d^0=0$.
All steps
use exact integral
finite matrix
operations.
\end{proof}

\begin{theorem}[Analytic integral Smith theorem for both odd sheets]
\label{thm:dyadic-f135-smith-9-27}
For each $n=9,27$,
the degree-one--two
matrix of
\eqref{eq:dyadic-f135-27sheet-fox-matrix}
has \emph{exact}
$2$-adic
Smith factors
\begin{equation}\label{eq:dyadic-f135-two-smith-factors}
\boxed{
\operatorname{SNF}_{\mathbf Z_2}(d_n^1)
=(1^{\,2n-1},2),
\quad
\operatorname{SNF}_{\mathbf Z_2}(d_n^0)
=(1^{\,n-1},0).
}
\end{equation}
Thus the finite
marked complex
has cohomology
\begin{equation}\label{eq:dyadic-f135-two-cohom-ledgers}
\boxed{
\begin{array}{c|ccc}
n&H^0&H^1&H^2\\ \hline
9&
\mathbf Z_2&
\mathbf Z_2^{10}&
\mathbf F_2\\
27&
\mathbf Z_2&
\mathbf Z_2^{28}&
\mathbf F_2.
\end{array}}
\end{equation}
These ledgers are
proved by \emph{integral
elementary operations},
not by relying only
on a finite-precision
matrix rank.

For both covers,
the source-owned
functional on the
full tame/wild
relation-face module
\begin{equation}\label{eq:dyadic-f135-odd-degree-top-trace}
\boxed{
\mathfrak t_n(u,v)
=\sum_{i=1}^{n}
\bigl(v_i+3u_i\bigr)
\pmod2,
\qquad
(u,v)\in V_n\oplus V_n,
}
\end{equation}
annihilates
$\operatorname{im}d_n^1$
and induces the
primitive top
cohomology isomorphism
\[
\mathfrak t_n:
H^2(\mathcal J_n)
\xrightarrow{\;\sim\;}\mathbf F_2.
\]
More precisely
\begin{equation}\label{eq:dyadic-f135-top-trace-integral-identity}
\boxed{\qquad
\sum_i
\Bigl((d_n^1)_{\rm wild}
+3(d_n^1)_{\rm tame}\Bigr)_i
=2\sum_i c_i
\qquad}
\end{equation}
for the four
degree-one inputs
$(a,b,c,d)$.
\end{theorem}

\begin{proof}
Write
\[
W_n=\operatorname{im}P_n
=V_n^{\langle\Upsilon_n\rangle},
\qquad
Q_n=\ker P_n.
\]
The order-nine
cyclic operator
$\Upsilon_n-I$
is an integral
automorphism of
$Q_n$:
the factor
$(T^9-1)/(T-1)$
has value $9$
at $T=1$,
an odd unit,
so its
augmentation
and complementary
summands split
over $\mathbf Z_2$.
The tame normalizer
$\Sigma_n$ preserves
$W_n$ and $Q_n$.
On $Q_n$, the
first tame Fox
block
$\Sigma_n^{-1}(\Upsilon_n-I)$
is invertible.
On $W_n$, the
second tame block
reduces to
$\Sigma_n^{-1}-2I$,
also invertible
because it is
congruent to
$\Sigma_n^{-1}$
modulo two.
Therefore the
tame differential
$V_n^2\to V_n$
is split surjective
with precisely $n$
unit pivots.

Restrict next to
its kernel.
The tame equation
forces the
$W_n$ component
of the $\tau$
cochain $b$ to be
zero.
On $Q_n$ the
wild $x_1$ block
$\Sigma_n^{-1}-P_n$
is $\Sigma_n^{-1}$,
hence supplies
$n-r$ unit
pivots, where
\[
r=\operatorname{rank}W_n
=\begin{cases}
1,&n=9,\\
3,&n=27.
\end{cases}
\]
On $W_n$ the
wild differential
now reduces to
\[
2c_W+
(\Sigma_n^{-1}-I)d_W.
\]
The $\Sigma_n$
action cyclically
permutes the
$r$ odd-inertia
orbits, and
$r$ is odd.
Thus
\[
W_n=\mathbf Z_2\mathbf1
\oplus W_n^{\rm aug},
\]
where
$W_n^{\rm aug}$
is the primitive
sum-zero
submodule.
The map
$\Sigma_n^{-1}-I$
is an integral
automorphism on
$W_n^{\rm aug}$,
giving $r-1$
further units.
On the remaining
constant line the
coefficient of
$c_W$ is exactly
$2$.
Altogether
$d_n^1$ has
\[
n+(n-r)+(r-1)=2n-1
\]
unit factors
and exactly
one factor two.
This proves its
Smith form.

The degree-zero
differential
$d_n^0$ has
$n-r$ unit
factors from
$\Upsilon_n-I$
on $Q_n$, and
$r-1$ from
$\Sigma_n-I$
on $W_n^{\rm aug}$.
Its kernel is
the constant
line
$\mathbf Z_2\mathbf1$,
so it has
$n-1$ unit factors
and one zero.
Because
$d_n^1d_n^0=0$,
the degree-one
cohomology is
torsion-free of
rank
$4n-2n-(n-1)
=n+1$.
The cokernel
of $d_n^1$
is $\mathbf F_2$,
giving both ledgers.

For the face trace,
every permutation
preserves the
sum of coordinates:
$\mathbf1^t\Sigma_n=
\mathbf1^t\Upsilon_n=
\mathbf1^t$,
and
$\mathbf1^tP_n=\mathbf1^t$.
Summing the tame
Fox relation gives
$-\sum_i b_i$;
summing the wild
relation gives
$3\sum_i b_i+
2\sum_i c_i$.
The stated
identity follows.
Reduction modulo
two annihilates
all top coboundaries.
It is primitive,
since it maps a
single wild
relation-face
basis vector to
one mod two.
Because the
top cokernel
has order two,
this proves
the claimed
isomorphism.
\end{proof}

\begin{theorem}[Actual pro-$2$ Schreier--Tietze reduction in nine and twenty-seven sheets]
\label{thm:dyadic-f135-group-fox-9-27}
For $H_n=G_{F_9}$
when $n=9$, and
$H_n=G_{K_{27}}$
when $n=27$,
take the \emph{actual}
marked
Reidemeister--Schreier
rewriting of both
full-group
Roe--Turturean
relations using
the coset representatives
of
Theorem~\ref{thm:dyadic-f135-permutation-9-27}.
After contracting
$n-1$ spanning-tree
edges, it is an
exact presentation
of the maximal
pro-$2$ quotient
$P_n=H_n(2)$
with
\[
3n+1
\quad\text{pro-$2$
generators, and}\quad
2n
\quad\text{marked
relation words}.
\]
These $2n$ actual
relators have
primitive augmented
Fox rank exactly
$2n-1$ over
$\mathbf F_2$.
The $n$ tame
relation directions
are primitive,
and modulo their
span, $n-1$ of
the wild directions
are primitive.
Therefore genuine
pro-$2$
Nielsen--Tietze
operations eliminate
\emph{exactly}
\[
n+(n-1)=2n-1
\]
relation/generator
pairs and leave a
minimal one-relator
presentation
\begin{equation}\label{eq:dyadic-f135-minimal-9-27-presentation}
\boxed{
\begin{aligned}
P_9&\cong
\langle z_1,\ldots,z_{11}
\mid\mathscr R_9\rangle_{\mathrm{pro}-2},\\
P_{27}&\cong
\langle z_1,\ldots,z_{29}
\mid\mathscr R_{27}\rangle_{\mathrm{pro}-2}.
\end{aligned}}
\end{equation}
The relators
$\mathscr R_9,\mathscr R_{27}$
are obtained by a
fully specified
\emph{finite-stage algorithm}:
Schreier word scan,
augmented mod-two
unit-minor pivot,
pro-$2$ Burnside
basis automorphism,
and inverse-limit
generator substitution.
They are not
claimed to be
closed finite
ordinary words.

For any finite
projective complete
coefficient module
$M$ whose
$H_n$ action has
pro-$2$ image,
there is a
\emph{chosen strict
integral Fox chain
equivalence}
\begin{equation}\label{eq:dyadic-f135-group-level-chain-split}
\boxed{
\begin{aligned}
C_{\rm cover}^\bullet(H_n,M)
\cong {}&
C_{\rm Fox}^\bullet(P_n,M)\\
&\oplus
[M^{n-1}\xrightarrow I M^{n-1}]_{[0,1]}\\
&\oplus
[M^{2n-1}\xrightarrow I M^{2n-1}]_{[1,2]}.
\end{aligned}}
\end{equation}
The right-hand
minimal Fox complex
is an exact
three-term
Demu\v{s}kin
cochain carrier.
This comparison
is an \emph{actual}
completed-group-algebra
Fox--Shapiro
comparison,
rather than a
numerical
cohomology agreement.
Every prescribed
finite pro-$2$
word or coefficient
jet is computed
by a terminating
sequence of
finite quotient
operations, as in
fixed133.

In particular
the genuine
twenty-seven-sheet
Sylow source
has degrees
\[
(27,108,54)
=(1,29,1)
+(26,26,0)
+(0,53,53).
\]
No functorial
preferred minimal
basis across
arbitrary dyadic
fields is inferred.
\end{theorem}

\begin{proof}
The usual finite-index
Reidemeister--Schreier
scan of the four
marked source
generators gives
$4n$ directed
edges, and the
two global relators
lift to $2n$
relation faces.
Contracting a
spanning tree on
the $n$ cosets
leaves one vertex
and $4n-(n-1)
=3n+1$ edges.

It remains to
check that the
marked wild
pro-$2$ condition
does not impose
extra unlisted
relations on the
pro-$2$ quotient.
For any finite
$2$-group quotient
$P_n\to Q$,
the induced
coset/wreath
representation of
the full
\emph{unmarked}
two-relator source
lands in
$Q^n\rtimes Q_9$.
The wild generators
have trivial
$Q_9$ permutation
image; their
normal closure
lies in the
normal finite
$2$-group
$Q^n$.
Hence this
finite representation
satisfies the
marked Roe--Turturean
pro-$2$ wild
condition.
Conversely every
actual $H_n$
quotient satisfies
the Schreier
relators.
These two finite
quotient functors
coincide, so after
pro-$2$ completion
the actual indexed
Schreier relators
present $P_n$.

Their augmented
Fox matrix over
$\mathbf F_2$
is the ordinary
two-cell
Schreier relation
matrix: equivalently
it is the
degree-one--two
coefficient
boundary of
the permutation
coinduced module
$V_n$ after
the $n-1$
tree-edge unit
contractions.
The latter matrix
is exactly
\eqref{eq:dyadic-f135-27sheet-fox-matrix}.
The Smith theorem
\ref{thm:dyadic-f135-smith-9-27}
shows that its
augmented mod-two
rank is $2n-1$,
including $n$
primitive tame
directions and
$n-1$ primitive
wild directions.
Thus an explicit
$(2n-1)\times(2n-1)$
unit minor exists.
Choose the
lexicographically
first such minor
as a deterministic
pivot convention.

The associated
$2n-1$
relator words,
together with
the remaining
$n+2$ generator
letters, are
a basis of the
Frattini quotient of
the free pro-$2$
Schreier group.
The pro-$2$
Burnside basis
theorem lifts
them to a
free pro-$2$
basis automorphism.
The inverse
automorphism
turns these
$2n-1$ relators
into generator
letters; killing
them is an
honest pro-$2$
Tietze elimination,
leaving precisely
one relator on
$n+2$ letters.
Since $P_n$ is
the local
Demu\v{s}kin group
from
Theorem~\ref{thm:dyadic-f135-actual-27-sylow-field},
the one remaining
relator is genuine
and defines its
minimal pro-$2$
resolution.

Each group-level
free-basis
automorphism
induces an
invertible
completed Fox
Jacobian over
$\mathbf Z_2[[P_n]]$.
For a
generator-letter
relator the
Fox boundary
has an identity
pivot; it is
cancelled by a
strict chain
row/column
operation.
The tree
contractions and
the $2n-1$
relation--edge
contractions give
exactly
\eqref{eq:dyadic-f135-group-level-chain-split}.
The remaining
Demu\v{s}kin
Fox complex
computes
$R\Gamma(P_n,M)$
by the pro-$2$
Poincar\'e
duality theorem
\cite{NSW,LabuteClassification}.

For any finite
source word
precision, the
Schreier scan
and each
$1/3^k$ pro-$2$
power are finite
operations.
The chosen
pro-$2$ basis
automorphism is
bijective on
every finite
characteristic
quotient, so
its inverse
can be computed
by finite search
or successive
Frattini lifts.
Completed Fox
differentiation
and the unit
pivot inverses
likewise terminate
at each prescribed
finite quotient
and coefficient
precision.
Their compatible
inverse limits
give the strict
integral comparison.
\end{proof}

\begin{theorem}[Four nested integral Hecke projectors on the real Sylow cover]
\label{thm:dyadic-f135-four-heeke-projectors}
On the degree-$27$
coset basis $e_{i,j}$
of
\eqref{eq:dyadic-f135-27-coset-representatives},
define the actual
$F_9$-coset label
\begin{equation}\label{eq:dyadic-f135-27-to-9-label}
\boxed{\qquad
k(i,j)=2^j i\bmod9.
\qquad}
\end{equation}
Its reduction modulo
three is the
$F_3$-coset label.
For $r=1,3,9$
define the
fiber-average
idempotent
\begin{equation}\label{eq:dyadic-f135-four-fiber-averages}
\boxed{
(Q_r)_{ab}
=
\begin{cases}
r/27,
&k(a)\equiv k(b)\pmod r,\\
0,&\text{otherwise},
\end{cases}
\qquad Q_{27}=I_{27}.
}
\end{equation}
All entries are
integral in
$\mathbf Z_2$,
and these
operators commute
with the full
$\Gamma$ action
on the coset
permutation module.
They satisfy
the \emph{strict}
nested identities
\begin{equation}\label{eq:dyadic-f135-nested-average-relations}
\boxed{\qquad
Q_rQ_s=Q_{\min(r,s)},
\qquad
\operatorname{rank}Q_r=r
\quad(r,s\in\{1,3,9,27\}).
\qquad}
\end{equation}
Their successive
orthogonal differences
\begin{equation}\label{eq:dyadic-f135-four-orthogonal-layers}
\boxed{
E_0=Q_1,\quad
E_1=Q_3-Q_1,\quad
E_2=Q_9-Q_3,\quad
E_3=I_{27}-Q_9
}
\end{equation}
are mutually
annihilating
strict idempotents
with
\begin{equation}\label{eq:dyadic-f135-primitive-layer-ranks}
\boxed{\quad
(\operatorname{rank}E_0,
\operatorname{rank}E_1,
\operatorname{rank}E_2,
\operatorname{rank}E_3)
=(1,2,6,18).
\quad}
\end{equation}

For \emph{every}
complete finite
projective
$\Gamma$-coefficient
system $B$
whose restriction
to $H_{27}$
has pro-$2$
image, tensoring
these operators
with $1_B$
gives commuting
idempotents in
\emph{every}
degree of the
marked coinduced
Fox cochain
complex:
\begin{equation}\label{eq:dyadic-f135-four-heeke-cochain-projectors}
\boxed{
\mathcal Q_r^0=Q_r\otimes1_B,\quad
\mathcal Q_r^1=Q_r^{\oplus4}\otimes1_B,\quad
\mathcal Q_r^2=Q_r^{\oplus2}\otimes1_B.
}
\end{equation}
Their differentials
satisfy
\[
\mathcal Q_r^1d^0
=d^0\mathcal Q_r^0,
\qquad
\mathcal Q_r^2d^1
=d^1\mathcal Q_r^1.
\]
Under the strict
Fox--Shapiro
equivalence of
Theorem~\ref{thm:dyadic-f135-group-fox-9-27},
the four projections
correspond to
the successive
odd-index
restriction--corestriction
idempotents for
\[
\Gamma\supset H_3
\supset H_9
\supset H_{27}.
\]
More precisely, on
$R\Gamma(H_{27},B)$
the projector $Q_r$,
where
$H_1:=\Gamma$,
represents
\begin{equation}\label{eq:dyadic-f135-derived-transfer-correspondence}
\boxed{
Q_r\simeq
\frac{r}{27}\,
\operatorname{res}_{H_{27}}^{H_r}
\operatorname{cor}_{H_{27}}^{H_r}.
}
\end{equation}
No division by
two is used.
The unnormalized
Hecke operators
$\mathcal A_r=(27/r)Q_r$
and
$\mathcal T_r=\mathcal A_r-I$
obey
\begin{equation}\label{eq:dyadic-f135-coherent-Hecke-polynomial}
\boxed{
\begin{gathered}
\mathcal A_r^2=(27/r)\mathcal A_r,\\
\mathcal T_r^2=
(27/r-2)\mathcal T_r+
(27/r-1)I,\\
\mathcal A_r\mathcal A_s
=(27/\max(r,s))\mathcal A_{\min(r,s)}.
\end{gathered}}
\end{equation}
These are
\emph{literal}
chain identities,
not merely
relations between
operators on
cohomology.
\end{theorem}

\begin{proof}
For the coset
representative
$\tau^i\sigma^j$,
moving $\sigma^j$
to the left gives
\[
\tau^i\sigma^j
=\sigma^j\tau^{2^j i}.
\]
Thus the image
of the degree-27
coset in
$H_9\backslash\Gamma$
has label
$k=2^ji\bmod9$.
Reduction modulo
three gives the
$H_3$ quotient.
The fibers for
$r=1,3,9,27$
have sizes
$27,9,3,1$.
Averaging over
a fiber defines
an integral
idempotent because
each size is odd.
The quotient maps
are $\Gamma$-equivariant;
therefore the fiber
averages commute
with every coset
permutation.
Averaging first
over a finer
fiber and then
over a coarser
fiber equals the
coarser average,
giving the
nested product
law and ranks.
Subtracting
successive
nested idempotents
gives the
orthogonal
differences
with the ranks
in
\eqref{eq:dyadic-f135-primitive-layer-ranks}.

The Fox complex
is functorial in
the coefficient
module.
Each $Q_r\otimes1_B$
commutes with
the action of
every marked
Galois generator,
hence with
every evaluated
completed Fox
relation derivative.
Applying it in
degrees zero,
one and two
proves all strict
cochain identities.
The coinduced
Fox--Shapiro
equivalence
identifies the
fiberwise unit
and sum maps
with restriction
and corestriction
for the corresponding
intermediate
field tower.
The normalized
fiber average
is their
composition
divided by
the odd index
$27/r$.
This proves
\eqref{eq:dyadic-f135-derived-transfer-correspondence}.
The Hecke
polynomials
and mixed
identities
follow by
substituting
$\mathcal A_r=(27/r)Q_r$
into the exact
nested
idempotent
relations.
\end{proof}

\begin{corollary}[Integral $H^1$ layers and top Tate-class concentration]
\label{cor:dyadic-f135-27sheet-cohom-layers}
For the tower
\[
\mathbf Q_2\subset
F_3\subset F_9\subset K_{27},
\]
the ordinary
rank-one integral
cohomology is
\[
\begin{array}{c|ccc}
K&H^0(G_K,\mathbf Z_2)&
H^1(G_K,\mathbf Z_2)&
H^2(G_K,\mathbf Z_2)\\ \hline
\mathbf Q_2&\mathbf Z_2&\mathbf Z_2^2&\mathbf F_2\\
F_3&\mathbf Z_2&\mathbf Z_2^4&\mathbf F_2\\
F_9&\mathbf Z_2&\mathbf Z_2^{10}&\mathbf F_2\\
K_{27}&\mathbf Z_2&\mathbf Z_2^{28}&\mathbf F_2.
\end{array}
\]
In the
strict decomposition
of
Theorem~\ref{thm:dyadic-f135-four-heeke-projectors},
all top-degree
two-primary torsion
is carried by
$E_0$, while
the three primitive
higher layers have
cohomology concentrated
in degree one,
of ranks
\begin{equation}\label{eq:dyadic-f135-layer-H1-ledger}
\boxed{\qquad
(2,\ 6,\ 18).
\qquad}
\end{equation}
Specifically, in
the derived category
of perfect
$\mathbf Z_2$
complexes,
\begin{equation}\label{eq:dyadic-f135-odd-tower-derived-splitting}
\boxed{
R\Gamma(H_{27},\mathbf Z_2)
\simeq
R\Gamma(\Gamma,\mathbf Z_2)
\oplus\mathbf Z_2^2[-1]
\oplus\mathbf Z_2^6[-1]
\oplus\mathbf Z_2^{18}[-1].
}
\end{equation}
The same
three primitive
degree-one ranks
hold for the
cyclotomic twist
$\mathbf Z_2(1)$.
Its top
integral Tate line
lies wholly in
the constant
coset component,
and its unique
degree-one
$\mathbf Z/2$
Kummer class
does as well.

For an explicit
Fox cochain with
trivial coefficients,
let $t,w\in V_{27}$
be the two
relation-face
vectors.
The primitive
top class
has the exact
source readout
\begin{equation}\label{eq:dyadic-f135-27face-Tate-readout}
\boxed{\qquad
[t,w]\longmapsto
\sum_i(w_i+3t_i)
\pmod2.
\qquad}
\end{equation}
Restriction of
the global marked
top class to
the twenty-seven
sheets multiplies
its source
sum by $27$
and thus acts as
the identity
on the surviving
$\mathbf F_2$
class.
\end{corollary}

\begin{proof}
Each field
has odd degree
over $\mathbf Q_2$,
so its
$2$-power roots
of unity are
exactly $\mu_2$.
Local class field
theory and
integral Tate
duality give
\[
H^0(G_K,\mathbf Z_2)=\mathbf Z_2,\quad
H^1(G_K,\mathbf Z_2)
=\mathbf Z_2^{[K:\mathbf Q_2]+1},\quad
H^2(G_K,\mathbf Z_2)=\mathbf F_2.
\]
The same ledgers
for the fields
$F_9,K_{27}$
were verified
at the finite
Fox level in
Theorem~\ref{thm:dyadic-f135-smith-9-27}.
Restriction is
split injective
at each odd-index
stage because
the normalized
corestriction
is its left inverse.
It is an
isomorphism
on the
one-dimensional
$\mathbf F_2$
top torsion,
and on $H^0$.
The successive
complementary
perfect complexes
therefore have
zero cohomology
except in degree
one, where their
free ranks are
$4-2=2$,
$10-4=6$,
$28-10=18$.
They are derived
equivalent to the
displayed free
degree-one modules.

For cyclotomic
coefficients,
local Kummer
theory gives
$H^0=0$,
$H^1=\mathbf Z_2^{d+1}
\oplus\mathbf F_2$,
and $H^2=\mathbf Z_2$.
Odd-index
restriction acts
invertibly on the
unique torsion
class and on
the top Tate line
(the latter by
multiplication by
the odd degree).
Thus all three
primitive
summands again
have cohomology
only in free
degree one,
with the same
rank differences.

The source
top functional
is
Theorem~\ref{thm:dyadic-f135-smith-9-27}.
A constant global
top face lifts
to the same
value on all
twenty-seven
indexed faces;
the top sum
therefore
multiplies by
$27$.
Reduction mod two
makes this the
identity on
$\mathbf F_2$,
and the
strict Hecke
projections
place the
top class in
the constant
summand.
\end{proof}

\begin{theorem}[A sufficient odd-index marked Schreier--Fox criterion]
\label{thm:dyadic-f135-general-primitive-schreier-criterion}
Let $H\subset G$
be an open
subgroup of
\emph{odd}
index $n$,
and suppose
$G$ has a
finite marked
profinite source
presentation,
with its
specified
normal-closure
conditions.
Assume that
\begin{enumerate}[label=\textnormal{(\roman*)}]
\item the marked
Reidemeister--Schreier
quotient on $H$,
after passing
to maximal
pro-$2$ targets,
has a genuine
finite pro-$2$
presentation with
$g$ generators
and $q$ relation
words;
\item its
augmented
Fox relation
matrix over
$\mathbf F_2$
has row rank
$q-1$, and
a chosen set
of $q-1$ primitive
relation words
gives a
unit Fox minor;
\item $H(2)$ is
a Demu\v{s}kin
group, and the
coefficient
modules under
consideration
have pro-$2$
$H$ action.
\end{enumerate}
Then a choice
of the primitive
minor determines
a finite-stage
computable
pro-$2$
Nielsen--Tietze
reduction to
a minimal
$(g-q+1)$-generator,
one-relator
Demu\v{s}kin
presentation.
The full marked
Schreier covering
Fox complex is
integrally
chain equivalent
to the genuine
minimal Fox
resolution plus
the selected
contractible
relation--edge
and tree-edge
pairs.
Moreover the
coset-average
\begin{equation}\label{eq:dyadic-f135-general-odd-index-projector}
\boxed{\qquad
e_{G,H}=
\frac1n\,\iota_{\rm const}
\circ\operatorname{sum}_{G/H}
\qquad}
\end{equation}
is a strict
integral
coefficient-level
chain idempotent
representing
$(1/n)\operatorname{res}
\operatorname{cor}$.
For a tower of
such marked
odd-index
covers, the
fiber-average
projectors are
strictly nested
and satisfy
the relations
of
\eqref{eq:dyadic-f135-nested-average-relations}
after identifying
their quotient
coset sets.

This is a
\emph{sufficient
criterion},
not a claim
that every
finite-index
marked presentation
automatically
has a primitive
Fox minor of
the required
rank.
It is verified
here for the
actual nine-
and twenty-seven-sheet
tame quotient
$Q_9$.
\end{theorem}

\begin{proof}
The primitive
row rank
condition permits
$\,q-1\,$ relator
words to be
substituted for
$\,q-1\,$ free
pro-$2$ generator
letters with
invertible
Frattini matrix.
The pro-$2$
Burnside basis
theorem makes
this a free
pro-$2$
automorphism.
Killing the
resulting
generator-letter
relators leaves
one relation
on $g-q+1$
generators;
minimality
follows from
the augmented
rank condition.
The Fox chain
rule gives
the corresponding
invertible
completed-group-algebra
chain changes,
and the eliminated
generator-letter
relations give
identity
contractible
blocks.
The remaining
Demu\v{s}kin
one-relator Fox
complex is exact
by its local
duality resolution.

The ordinary
sum and constant
unit maps on
the permutation
module are
$G$-equivariant.
Since $n$ is
a $2$-adic unit,
their normalized
composite is
an integral
idempotent,
and Fox
functoriality
makes it a
strict chain map.
Under Shapiro,
the constant
unit map is
restriction and
the unnormalized
sum is
corestriction.
For a tower,
averaging over
a fine fiber
followed by
averaging over
a coarse one
equals the
coarse average.
This proves
the stated
nested relations.
\end{proof}

\begin{theorem}[The full three-power family of odd-index Sylow--Fox towers]
\label{thm:dyadic-f135-all-three-power-sylow-towers}
For every integer
$k\ge1$, choose
$\alpha_k^{\,3^k}=2$
and set
\begin{equation}\label{eq:dyadic-f135-uniform-three-power-fields}
\boxed{
\begin{gathered}
F_k=\mathbf Q_2(\alpha_k),\\
U_{3^{k-1}}=
\text{the unramified degree-}3^{k-1}
\text{ extension of }\mathbf Q_2,\\
K_k=F_kU_{3^{k-1}},\qquad
n_k=[K_k:\mathbf Q_2]=3^{2k-1}.
\end{gathered}}
\end{equation}
Then the Galois
closure
$L_k=F_k(\zeta_{3^k})$
has tame quotient
\begin{equation}\label{eq:dyadic-f135-uniform-tame-quotient}
\boxed{
\operatorname{Gal}(L_k/\mathbf Q_2)
\simeq C_{3^k}\rtimes
C_{2\cdot3^{k-1}},
\qquad
\sigma^{-1}\tau\sigma=\tau^2.
}
\end{equation}
The subgroup
$\langle\sigma^{3^{k-1}}\rangle$
is a Sylow $2$-subgroup;
its fixed field is
precisely $K_k$.
Thus $G_{K_k}$
is the \emph{true}
Sylow preimage
of index $n_k$,
not merely a
chosen odd-degree
radical subgroup.

Let $V_k=\mathbf Z_2^{n_k}$
be the coinduced
trivial permutation
lattice.
Its tame inertia
$\Upsilon_k$
has exactly
$r_k=3^{k-1}$
orbits, each of
length $3^k$,
and $\Sigma_k$
permutes those
orbits cyclically.
The normalized
odd-inertia projector
\[
P_k=\frac1{3^k}
\sum_{a=0}^{3^k-1}\Upsilon_k^a
\]
is integral
of rank $r_k$.
For the true
marked coinduced
two-relator
Fox differential,
one has the
\emph{same}
symbolic formula
\eqref{eq:dyadic-f135-27sheet-fox-matrix}
with
$(n,\Sigma_n,\Upsilon_n,P_n)$
replaced by
$(n_k,\Sigma_k,\Upsilon_k,P_k)$.
The entire
Smith calculation
therefore extends
without a new
finite matrix
search:
\begin{equation}\label{eq:dyadic-f135-uniform-odd-Smith}
\boxed{
\operatorname{SNF}(d_k^1)
=(1^{2n_k-1},2),
\qquad
\operatorname{SNF}(d_k^0)
=(1^{n_k-1},0).
}
\end{equation}
The actual
marked group
Schreier relators
have mod-two
primitive row
rank $2n_k-1$,
so the
pro-$2$ Burnside
criterion gives
a finite-stage
computable
\emph{one-relator}
presentation of
$G_{K_k}(2)$
with $n_k+2$
minimal generators,
and its genuine
integral
Demu\v{s}kin
Fox chain
resolution after
$n_k-1$
tree contractions
and $2n_k-1$
relation cancellations.

There is a
chain of
$2k-1$ nested
odd-index-three
fields
\begin{equation}\label{eq:dyadic-f135-uniform-field-chain}
\mathbf Q_2
\subset F_1\subset\cdots\subset F_k
\subset F_kU_3\subset\cdots
\subset F_kU_{3^{k-1}}=K_k,
\end{equation}
where the
second segment
is empty for
$k=1$.
The successive
normalized
coset-average
idempotents
are strictly
commuting cochain
endomorphisms,
and their
orthogonal layers
have exact
integral ranks
\begin{equation}\label{eq:dyadic-f135-uniform-hecke-ranks}
\boxed{\qquad
1,\ 2,\ 6,\ 18,\ \ldots,\
2\cdot3^{2k-2},
\qquad
1+\sum_{j=1}^{2k-1}
2\cdot3^{j-1}
=n_k.
\qquad}
\end{equation}
For trivial
integral coefficients
the nonconstant
layers have
cohomology in
degree one only,
with free ranks
$2\cdot3^{j-1}$,
while the unique
top $\mathbf F_2$
class is carried
by the constant
coset line.
The same free
primitive-layer
ranks hold for
cyclotomic coefficients.
The coefficient-uniform
coinduced
chain completeness
holds on any
framed residual
sector whose
restriction to
the Sylow
preimage $G_{K_k}$
has pro-$2$ image.

This is an
\emph{infinite,
explicit odd-index
family} within
the same marked
tame source,
not a
classification
of all finite
dyadic extensions.
Its minimal
Fox bases remain
choice-dependent
and finite-jet
computable.
\end{theorem}

\begin{proof}
The polynomial
$X^{3^k}-2$
is Eisenstein
and gives a
totally tame
extension of
degree $3^k$.
The cyclotomic
factor is
unramified,
of degree
\[
\operatorname{ord}_{3^k}(2)
=2\cdot3^{k-1},
\]
by the elementary
lifting-the-exponent
formula for
$2^{2\cdot3^j}-1$.
The Galois
closure therefore
has the displayed
semidirect product.
The unique
order-two subgroup
of its chosen
Frobenius
complement is
$\langle\sigma^{3^{k-1}}\rangle$.
It fixes the
chosen radical
root and
the unramified
degree-$3^{k-1}$
subfield, so
its fixed field
is $K_k$ and
its index is
$3^k3^{k-1}
=3^{2k-1}$.

The inertia
subgroup
$C_{3^k}$
acts freely
in each orbit
of the
$G_{K_k}\backslash\Gamma$
coset set.
There are
$r_k=3^{k-1}$
orbits, and
the remaining
Frobenius quotient
permutes them
as a single
cyclic orbit.
Because the
residual Frobenius
has order
$2\cdot3^{k-1}$,
its pro-$2$
component has
linear order
two.
Thus the affine
marked-word
proof of
Theorem~\ref{thm:dyadic-f135-27sheet-fox-block-formula}
uses precisely
the same
identities:
the pro-$2$
inertia projection
is the normalized
$3^k$-cycle norm,
the square of
the pro-$2$
Frobenius
component is
a pure translation,
and its conjugation
of a wild translation
is trivial.
This gives
the full
symbolic Fox
matrix for
every $k$.

The integral
Smith proof of
Theorem~\ref{thm:dyadic-f135-smith-9-27}
only uses the
odd-order inertia
projector,
the invertibility
of $\Upsilon-I$
on its complement,
the invertibility
of $\Sigma^{-1}-2I$
on its invariant
part, and the
cyclic permutation
of an
\emph{odd} number
$r_k$ of inertia
orbits.
All four
properties remain
true, and the
same calculation
gives $n_k$
tame units,
$n_k-r_k$
wild units,
$r_k-1$
additional
wild units,
and a final
factor two.
The degree-zero
image has
$n_k-1$
primitive
unit directions.
This proves
\eqref{eq:dyadic-f135-uniform-odd-Smith}.

For the genuine
pro-$2$ group
presentation,
the marked
finite-wreath
argument from
Theorem~\ref{thm:dyadic-f135-group-fox-9-27}
works unchanged.
The odd tame
element $\tau^{3^k}$
must vanish
in every finite
pro-$2$ quotient
of the Sylow
preimage:
its conjugate
by the Sylow
element
$\sigma^{3^{k-1}}$
is its
$2^{3^{k-1}}$th
power, an
even power.
A finite
$2$-group
element cannot
be conjugate
to a positive
even power
unless it is
trivial.
Therefore
the tame pro-odd
source condition
causes no
additional
pro-$2$
relation.
All induced
wild images
belong to the
pro-$2$ base
of the finite
wreath quotient,
so the genuine
Schreier relation
functor is
as stated.
The augmented
Fox rank then
gives the
same primitive
pro-$2$
Nielsen--Tietze
elimination.
Since $K_k$
has odd
degree,
its cyclotomic
orientation
is full and
$q(K_k)=2$,
while its
Demu\v{s}kin
generator rank
is $n_k+2$;
the remaining
one-relator
Fox resolution
is exact.

Finally the
field chain
\eqref{eq:dyadic-f135-uniform-field-chain}
has successive
degrees three.
Each intermediate
quotient of
the full coset
set is
$\Gamma$-equivariant.
Averaging along
its fibers
is integral,
idempotent,
and strictly
compatible
with all
completed Fox
differentials.
Their nested
differences
have ranks
$3^j-3^{j-1}
=2\cdot3^{j-1}$,
giving
\eqref{eq:dyadic-f135-uniform-hecke-ranks}.
Odd-index
restriction and
normalized
corestriction
split
cohomology
at every stage.
Local
class field
theory,
Kummer theory
and integral
Tate duality
show that the
degree-zero and
top classes
occur only
in the constant
summand;
the positive
primitive layers
have the stated
free degree-one
ranks.
The complete
coefficient
comparison follows
from the actual
pro-$2$
Schreier--Fox
resolution and
strict constant-coset
projection, as in
Theorem~\ref{thm:dyadic-f135-four-heeke-projectors}.
\end{proof}

\begin{warning}[Odd degree is not automatically a Sylow preimage]
\label{warn:dyadic-f135-nine-versus-27}
The degree-nine
field $F_9$
has odd index
but its
residual tame
stabilizer $C_6$
contains an
odd-order
subgroup.
Thus for a
general framed
residual $Q_9$
representation
its restricted
$H_9$ image
need \emph{not}
be pro-$2$.
The strict
minimal
pro-$2$ Fox
comparison on
$H_9$
is valid only
for coefficients
that actually
factor through
$H_9(2)$.
By contrast
$H_{27}$ is
the \emph{true}
Sylow-$2$ preimage
of $Q_9$;
every complete
framed
$Q_9$ residual
coefficient
representation
restricts to
a pro-$2$
image on
$H_{27}$.
The twenty-seven-sheet
comparison therefore
applies throughout
those residual
coefficient sectors.

The new
group-level
reduction is
choice-dependent
at its Schreier
and Nielsen
stages, and
the higher
overlap
coherences
between
different
odd-index
towers are
not asserted
to possess one
preferred
closed
word-level
formula.
The
strictly nested
coefficient
Hecke idempotents
are intrinsic
to the
chosen tower
of subgroups;
they do not
by themselves
produce a
canonical
minimal Fox
basis for
every dyadic
extension.
\end{warning}
